\documentclass[11pt]{article}
\addtolength{\textwidth}{4cm}
\addtolength{\textheight}{5cm}
\addtolength{\oddsidemargin}{-2cm}
\addtolength{\evensidemargin}{-4cm}
\addtolength{\topmargin}{-2cm}

%\setcounter{tocdepth}{2}
%\renewcommand{\baselinestretch}{2}

\begin{document}

Computer algebra tools have been used since a long time
for theoretical studies of parallel manipulators like computing
the maximum number of solutions for the direct kinematic problem.
Recent progress in algorithms for computing Gr\"oebner basis
and for the exact study of real roots have enhanced the computational
capabilities of dedicated computer algebra systems for the study
of the direct kinematics problem, so that now we can start using such
tools in practice for studying real and industrial situations.

The first part will be devoted to the description of
new algorithms and software solutions for the study of the direct
kinematics problem .
We will show how we can simplify efficiently the related systems of
algebraic equations and then isolate all their real roots in an exact
way. The first step is equivalent to eliminate variables of the
 the algebraic system that modelises any direct kinematics problem
 $\{h_1(X_1,\ldots , X_9)=,\cdots =h_9(X_1,\ldots , X_9)=0\}$.


\begin{itemize}
\item the {\bf elimination phase} produces an equivalent system
(Rational Univariate Representation) in the form :
$$
\left \{
\begin{array}{l}
f(T)=0 \\
X_1=g_1(T)/g(T)\\
\vdots \\
X_9=g_9(T)/g(T)
\end{array}
\right .
$$

\item the {\bf isolation phase} takes the result of the elimination
phase , isolates the real roots (with multi-precision rational bounds)
of the first polynomial ($f$) and computes isolation intervals for all
the coordinates $X_1,\ldots,X_9$, providing in addition an algorithm
for refining as precisely as we want the solutions.
\end{itemize}

These algorithms are implemented in FGb/RS software and it takes few
seconds to compute all the real roots for practical examples.

The second part will describe an extension of such a computation to path
planning. In particular we will demonstrate how the elimination phase
of the computation can be used safely for choosing the right path (among
the 40 possibilities) given a parametric description of any trajectory
and an initial position.

The third part will be devoted to the description of an industrial example

on which such a algorithms have been applied, providing at the end a kind
of
classification of the capabilities of the given parallel manipulator
depending on the available hardware and real-time numerical methods.

\end{document}
